% PyPSA-Eur run configuration for the nodal baseline (Section \ref{sec:Baseline}).
% Values taken from results/de-nodal-2025-hourly/configs/config.base_s_30_elec_.yaml
\begin{table}[t]
  \centering
  \caption{Configuration of the PyPSA-Eur baseline runs. Both runs share all
  settings except the spatial clustering. Note the deliberate split between the
  \emph{system year} 2025 (capacities, costs, CO\textsubscript{2} budget) and the
  \emph{weather year} 2013 (renewable and load time series); see the discussion in
  Section~\ref{sec:Baseline}.}
  \label{tab:pypsa-settings}
  \begin{tabular}{@{}ll@{}}
    \toprule
    \textbf{Setting} & \textbf{Value} \\
    \midrule
    PyPSA-Eur version        & v2026.02.0 \\
    Foresight                & overnight (single horizon) \\
    System / price year      & 2025 (\texttt{planning\_horizons}) \\
    Weather year             & 2013 \\
    Snapshots                & 2013-01-01 to 2014-01-01, hourly \\
                             & (8760\,h, leap day dropped) \\
    Atlite cutout            & \texttt{europe-2013-sarah3-era5} \\
    Countries                & DE + FR, BE, NL, LU, CH, \\
                             & AT, CZ, PL, DK (10 total) \\
    Spatial clustering       & 20 and 30 nodes \\
                             & ($\rightarrow$ 10 and 17 DE nodes) \\
    Transmission limit (\texttt{ll})    & v1.0 (no grid expansion) \\
    Transmission losses      & linearised, 2 tranches \\
    CO\textsubscript{2} budget (2025)   & 0.648 rel.\ to 1990 \\
    Load shedding            & enabled, 100\,000\,\euro/MWh \\
    Solver                   & Gurobi, barrier (method 2) \\
    Random seed              & 123 \\
    \bottomrule
  \end{tabular}
\end{table}
